\documentclass[a4paper,12pt]{ctexart}
\usepackage{../../teaching-style}

\begin{document}
\pagestyle{empty}

\maintitle{统计与概率初步}
\begin{center}
  \large 低学段
\end{center}
\vspace{0.3cm}

\chaptertitle{第一章\quad 数据的收集与整理}

\sectiontitle{1.1\quad 数据的收集}

\knowbox{
\textbf{调查方式}\\
全面调查（普查）：对所有对象逐一调查。\\
抽样调查：从总体中抽取一部分样本进行调查。\\
\textbf{总体}：调查对象的全体。\\
\textbf{样本}：从总体中抽取的一部分个体。\\
\textbf{个体}：总体中的每个对象。}

\sectiontitle{1.2\quad 数据的整理}

\knowbox{
用表格整理数据：列出分类、划记、频数。\\
\textbf{频数}：每个对象出现的次数。\\
\textbf{频率}：频数$\div$总数。\\
频数之和$=$总数，频率之和$=1$。}

\sectiontitle{习题1}

\begin{exercise}
\item 调查班级同学最喜欢的水果，设计一个调查表并填写。
\item 下列调查适合用普查还是抽样调查？\\
  (a) 检查灯泡的使用寿命；\\
  (b) 人口普查；\\
  (c) 检查某批食品的合格率。
\item 某班$40$人，统计喜欢体育项目的频数：足球$12$人，篮球$10$人，乒乓球$8$人，其他$10$人。求各项的频率。
\end{exercise}

\newpage

\chaptertitle{第二章\quad 统计图表}

\sectiontitle{2.1\quad 条形图}

\knowbox{用高度不同的直条表示数据的大小，直观比较各类数据。\\ 横轴表示类别，纵轴表示数量。}

\sectiontitle{2.2\quad 折线图}

\knowbox{用点连接成折线，表示数据随时间或其他顺序的变化趋势。}

\sectiontitle{2.3\quad 扇形图}

\knowbox{用扇形面积表示各部分占总体的百分比。\\ 圆心角 = $360^\circ\times$百分比。}

\sectiontitle{习题2}

\begin{exercise}
\item 某班数学成绩：$90$分以上$8$人，$80$~$89$分$15$人，$70$~$79$分$10$人，$60$~$69$分$5$人，$60$分以下$2$人。绘制条形图和扇形图（描述）。
\item 某地$1$~$6$月平均气温：$5,8,12,18,23,28$（$^\circ$C）。绘制折线图（描述），哪个月到哪个月升温最快？
\item 根据上题的数据，每季度的平均气温是多少？
\end{exercise}

\newpage

\chaptertitle{第三章\quad 数据的代表}

\sectiontitle{3.1\quad 平均数}

\knowbox{
\textbf{算术平均数}$\bar x=\dfrac{x_1+x_2+\cdots+x_n}{n}$\\
\textbf{加权平均数}：若各数有不同权重，$\bar x=\dfrac{w_1x_1+w_2x_2+\cdots+w_nx_n}{w_1+w_2+\cdots+w_n}$}

\sectiontitle{3.2\quad 中位数与众数}

\knowbox{
\textbf{中位数}：将数据从小到大排列，最中间的数（奇数个）或中间两个数的平均数（偶数个）。\\
\textbf{众数}：出现次数最多的数（可能不止一个）。\\
\textbf{区别}：平均数易受极端值影响，中位数和众数更稳定。}

\sectiontitle{习题3}

\begin{exercise}
\item 求$5,8,12,15,20$的平均数、中位数和众数。
\item 某组$10$个数据：$3,5,5,6,7,7,7,8,9,10$，求它们的平均数、中位数和众数。
\item 在$5$次测验中，小明得分$85,90,92,88,95$，求平均分。
\item 一组数据的众数是$8$，中位数是$7$，平均数是$7.5$，请写出一组满足条件的数据。
\item 老师说全班的平均分是$78$分，小明考了$85$分却低于平均分，可能吗？说明理由。
\end{exercise}

\newpage

\chaptertitle{第四章\quad 方差与标准差}

\sectiontitle{4.1\quad 方差}

\knowbox{
方差$s^2=\dfrac{(x_1-\bar x)^2+(x_2-\bar x)^2+\cdots+(x_n-\bar x)^2}{n}$\\
方差越大，数据越分散（波动越大）；方差越小，数据越集中（越稳定）。\\
\textbf{标准差}$s=\sqrt{s^2}$，与原始数据单位相同。}

\sectiontitle{习题4}

\begin{exercise}
\item 两组数据：$A: 5,5,5,5,5$；$B: 1,3,5,7,9$。求它们的方差。
\item 甲乙两名射击运动员，$10$次成绩：\\
  甲：$8,9,9,10,8,9,10,7,10,8$\\
  乙：$6,10,8,9,10,7,9,8,10,7$\\
  分别计算平均数和方差，谁的发挥更稳定？
\item 数据$a-3,a-1,a,a+1,a+3$的方差是多少？
\end{exercise}

\newpage

\chaptertitle{第五章\quad 概率初步}

\sectiontitle{5.1\quad 确定事件与随机事件}

\knowbox{
\textbf{必然事件}：一定会发生的事件（概率为$1$）。\\
\textbf{不可能事件}：一定不会发生的事件（概率为$0$）。\\
\textbf{随机事件}：可能发生也可能不发生的事件（概率在$0$到$1$之间）。}

\sectiontitle{5.2\quad 古典概型}

\knowbox{
如果一个随机试验满足：\\
1. 所有可能结果有限个；\\
2. 每个结果出现的可能性相等。\\
则事件$A$的概率$P(A)=\dfrac{\text{事件}A\text{包含的结果数}}{\text{所有可能结果总数}}$。}

\sectiontitle{5.3\quad 概率树形图与列表法}

\knowbox{
\textbf{树形图}：分步列出所有可能的结果，适合两步或三步的试验。\\
\textbf{列表法}：用表格列出所有可能结果，适合两步试验。}

\sectiontitle{习题5}

\begin{exercise}
\item 判断下列事件是必然事件、不可能事件还是随机事件：\\
  (a) 太阳从东方升起；\\
  (b) 明天会下雨；\\
  (c) 掷骰子出现$7$点。
\item 一个口袋装有$2$个红球和$3$个白球，从中任取$1$个球，取到红球的概率是多少？
\item 同时掷两枚硬币，两枚都是正面朝上的概率是多少？（用树形图或列表法）
\item 一个箱子中有$3$张卡片，分别写有$1$、$2$、$3$。从中先后抽出两张（不放回），求抽到的两张卡片数字之和为偶数的概率。
\item 在$1$到$10$中任选一个数，是质数的概率是多少？
\item 甲乙两人玩"石头剪刀布"游戏，甲获胜的概率是多少？
\end{exercise}

\newpage

\chaptertitle{第六章\quad 频率与概率}

\sectiontitle{6.1\quad 频率的稳定性}

\knowbox{
在大量重复试验中，事件发生的频率会稳定在某个常数附近，这个常数就是概率。\\
试验次数越多，频率越接近概率。\\
\textbf{大数定律直观理解}：掷硬币次数越多，正面朝上的频率越接近$\dfrac12$。}

\sectiontitle{6.2\quad 统计决策初步}

\knowbox{
利用统计数据和概率进行简单决策：\\
比较概率大小、分析数据趋势、做出合理选择。}

\sectiontitle{习题6}

\begin{exercise}
\item 小明掷一枚硬币$100$次，正面朝上$54$次。正面朝上的频率是多少？如果掷$1000$次，频率会怎样变化？
\item 某商场抽奖活动，共有$1000$张奖券，其中一等奖$1$名、二等奖$10$名、三等奖$50$名。\\
  (a) 中奖的概率是多少？\\
  (b) 中一等奖的概率是多少？
\item 甲乙两个水果店，甲店苹果合格率$95\%$，乙店$90\%$。你会在哪家买？为什么？
\item 天气预报明天下雨的概率是$80\%$，你会带伞吗？说说你的决策理由。
\end{exercise}

\end{document}
